Batteries in electric vehicles and grid storage are large-format cells, yet almost
all machine-learning work on battery health is built on small cylindrical cells
cycled in the laboratory, so their transfer to large-format cells is untested. To
close this gap, we release an 18-month dataset of 20 large-format LiFePO$_4$
prismatic cells (50\,Ah) cycled under three operating conditions. The data expose
the core difficulty: manufacturing variability, not temperature or model choice,
dominates how these cells age. Ten nominally identical cells reach final SOH values
that differ by 13.2$\times$ (7.7$\times$ even at the same cycle count), an order of
magnitude wider than any large-format LFP dataset we are aware of; across six
deep-learning models, the individual cell explains 25\% to 58\% of the prediction
error, the model only about 15\%. Can a new cell's health be predicted early,
without months of its own data? Our hybrid model, pretrained on many past cells,
reads a new cell's first 30 cycles through a fingerprint encoder and predicts its
SOH trajectory with no cell-specific training (RMSE 0.025, 38\% better than
training from scratch); fine-tuning raises this to 53.5\% over scratch on all 17
test cells, especially the fast-failing cells that drive warranty risk. It reports
calibrated confidence intervals (94\% coverage, but only
when training and deployment chemistries match), flags defective cells about 50
cycles early (0.78 AUC), and forecasts remaining life to a median of 49 cycles.
Because one-step accuracy saturates, even a linear baseline matches the best, we
judge models by these deployment capabilities and release a reproducible pipeline.
